\SourceLeaf{242}%
 \(x,y,z\) as functions of the indeterminates \(p,q\). The general expression for the magnitude of any line element suffices. Let us go on to some applications of this most important theorem.

Suppose that our curved surface can be developed onto another surface, curved or plane, so that to each point on the first surface determined by the coordinates \(x,y,z\) there corresponds a definite point on the second surface, with coordinates \(x',y',z'\). Evidently \(x',y',z'\) can also be regarded as functions of the indeterminates \(p,q\). For the element \(\BareRoot{(\mathrm{d}x'^2+\mathrm{d}y'^2+\mathrm{d}z'^2)}\) we thus obtain an expression of the form
\[
\BareRoot{(E'\mathrm{d}p^2+2F'\mathrm{d}p\mathord{.}\mathrm{d}q+G'\mathrm{d}q^2)}
\]
where \(E',F',G'\) likewise denote functions of \(p,q\). The very notion of the \textit{development} of one surface onto another makes it clear that corresponding elements on the two surfaces must be equal. We therefore have the identities
\[
E=E',\quad F=F',\quad G=G'
\]
Thus the formula of the preceding article leads directly to the remarkable

\textsc{Theorem.} \textit{If a curved surface is developed onto any other surface, the measure of curvature at each point remains unchanged.}

It is also evident that \textit{any finite part of the curved surface retains the same integral curvature after development onto another surface.}

Surfaces that can be developed onto a plane make up the special case to which geometers have so far confined their investigations. Our theory shows directly that the measure of curvature of such surfaces is \(=0\) at every point. Consequently, if their nature is described by the third method, the following holds everywhere
\[
\frac{\mathrm{d}\mathrm{d}z}{\mathrm{d}x^2}\mathord{.}\frac{\mathrm{d}\mathrm{d}z}{\mathrm{d}y^2}-(\frac{\mathrm{d}\mathrm{d}z}{\mathrm{d}x\mathord{.}\mathrm{d}y})^2=0
\]
This criterion, though long known, is generally proved with less rigor than one might wish, at least in our judgment.
\Article{13}
What we set out in the preceding article is connected with a particular way of considering surfaces that amply deserves careful study by geometers. Specifically, when a surface is viewed apart from its role as the boundary of a solid and considered as a
